% Part: first-order-logic
% Chapter: model-theory
% Section: substructures

\documentclass[../../../include/open-logic-section]{subfiles}

\begin{document}

\olfileid{mod}{bas}{sub}
\olsection{Sub\printtoken{p}{structure} (deelstructuren)}

Het !!{domain} van !!a{structure}~$\Struct{M}$ kan een deelverzameling
zijn van dat van een andere~$\Struct{M'}$.  Maar we beschouwen
$\Struct{M}$ vanzelfsprekend alleen als een ``deel'' van $\Struct{M'}$
als niet alleen $\Domain{M} \subseteq \Domain{M'}$ geldt, maar
$\Struct{M}$ en $\Struct{M'}$ ook ten minste op het gedeelde
deel~$\Domain{M}$ ``overeenkomen'' in de wijze waarop zij de symbolen
van de taal interpreteren.

\begin{defn}
\ollabel{defn:substructure}
Gegeven !!{structure}s $\Struct M$ en $\Struct M'$ voor dezelfde
taal~$\Lang L$ noemen we $\Struct M$ een \emph{sub!!{structure}} van
$\Struct M'$ en $\Struct M'$ een \emph{uitbreiding} van $\Struct M$,
genoteerd als $\Struct M \substruct \Struct M'$, precies dan als
\begin{enumerate}
\item $\Domain{M} \subseteq \Domain{M'}$,
\item voor elke constante $c \in \Lang L$ geldt $\Assign{c}{M} =
    \Assign{c}{M'}$;
\item voor elke $n$-plaatsige !!{function} $f \in \Lang L$ geldt
  $\Assign{f}{M}(a_1, \dots, a_n) = \Assign{f}{M'}(a_1, \dots, a_n)$
  voor alle $a_1$, \dots, $a_n \in \Domain{M}$.
\item voor elk $n$-plaatsig !!{predicate} $R \in \Lang L$ geldt $\langle
  a_1, \dots, a_n\rangle \in \Assign{R}{M}$ iff $\langle a_1, \dots,
  a_n\rangle \in \Assign{R}{M'}$ voor alle $a_1$, \dots, $a_n \in
  \Domain{M}$.
\end{enumerate}
\end{defn}

\begin{rem}
\ollabel{rem:substructure}
Als de taal geen constanten of !!{function}s bevat, bepaalt elke $N
\subseteq \Domain{M}$ een sub!!{structure}~$\Struct{N}$ van
$\Struct M$ met !!{domain}~$\Domain{N} = N$ door $\Assign{R}{N} =
\Assign{R}{M} \cap N^n$ te stellen.
\end{rem}

% prove something about this? Examples?

\end{document}
